\chapter*{第一部分导读：把计算环境变成科研工作流}
\addcontentsline{toc}{chapter}{第一部分导读：把计算环境变成科研工作流}

从这一部分开始，本书进入正文主线。Part I 的任务不是把 Linux、Git、Jupyter、FITS 和绘图各讲成一份工具说明书，而是帮助你建立一条最小科研计算工作流：知道文件在哪里，知道代码如何运行，知道数据如何进入程序，知道结果如何被记录和展示。

很多初学者会把科研计算理解成“学会几条命令”或“会写几段 Python”。这当然重要，但还不够。真正进入课程项目或科研训练时，更关键的问题常常是：

\begin{itemize}
    \item 当前目录是否就是项目目录；
    \item 这个结果来自哪个版本的代码和数据；
    \item notebook 里的探索代码能否整理成脚本；
    \item 文件读入之后是否真的满足字段、单位和元数据要求；
    \item 图表是否能支持一个清楚、有限、可复查的科学论断。
\end{itemize}

这些问题看起来琐碎，却决定了后面所有机器学习和科学解释是否可信。一个模型分数再漂亮，如果它来自混乱目录、未知数据版本、不可复跑 notebook 或错误单位，就不应被当作可靠结果。

\section*{本部分的学习顺序}

Part I 的五章可以看作一条从工作环境到科学表达的递进链：

\begin{enumerate}
    \item \textbf{Unix/Linux 与科研计算环境}：先建立项目目录、命令行、安全批处理、远程运行和日志意识；
    \item \textbf{Git 与可复现科研}：把代码、报告、notebook 和结果说明接入版本历史；
    \item \textbf{库、脚本与 Jupyter}：区分探索、脚本化和模块化，避免把临时 notebook 当成最终工作流；
    \item \textbf{文件 I/O、FITS 与科学数据入口}：从“文件能打开”推进到“数据结构、字段、单位和元数据可解释”；
    \item \textbf{可视化与科研图表}：把数值输出组织成图表证据，再从图表回到科学论断。
\end{enumerate}

这五章不要求你一下子成为系统管理员、软件工程师或专业数据工程师。它们的目标更具体：让你能独立维护一个小型科研项目，并让别人有可能复查你做过什么。

\section*{一个贯穿 Part I 的最小项目循环}

你可以把 Part I 的技能合并成下面这条循环：

\begin{enumerate}
    \item 用 Linux 命令建立清楚的项目目录；
    \item 用 Git 初始化版本历史，并写下第一份 README；
    \item 在 notebook 中探索数据和图表；
    \item 把稳定逻辑抽成脚本或函数；
    \item 读取数据时检查字段、单位、缺失值和元数据；
    \item 生成图表时保存代码入口、参数、输出路径和图表说明；
    \item 在一次清楚的 Git 提交中记录这个阶段的结果。
\end{enumerate}

这条循环本身就是后面所有案例的雏形。Part II 处理星表、图像、光谱、时间序列和实验数据时，仍然会回到它；Part III 训练机器学习模型时，也仍然要依赖它。

\begin{protip}[把工具学成流程，而不是清单]
本部分每章都会出现一些命令、代码片段或检查清单。不要把它们当作孤立知识点背诵。更好的学习方式是问：这一步在科研工作流里保护了什么？它是在保护数据、代码、环境、结果，还是科学解释？
\end{protip}

\section*{进入下一部分前的自检}

读完 Part I 后，你至少应该能够独立完成下面的小任务：

\begin{itemize}
    \item 在命令行里创建一个项目目录，并区分原始数据、脚本、notebook、日志和结果；
    \item 使用 \texttt{git status}、\texttt{git diff} 和 \texttt{git commit} 记录一次清楚的小改动；
    \item 解释 notebook、script 和 module 在科研项目中的分工；
    \item 读入一个 CSV 或教学版 FITS/cutout 数据，并写下最小数据检查记录；
    \item 画出一张带变量、单位、标题、图例和文字解释的科学图表；
    \item 用一两句话说明图表支持什么论断、证据在哪里、限制是什么。
\end{itemize}

如果这些任务还很吃力，不必急着进入复杂模型。先把 Part I 当作可复现实战的地基打牢，会比过早追逐高级算法更稳。
